\chapter{Andrei Okounkov}\label{ch2006andreiokounkov:andrei-okounkov-fields-medal-2006-a-comprehensive-report-on-his-work-and-impact}

\begin{center}\rule{0.5\linewidth}{0.5pt}\end{center}

\textbf{Prepared:} July 2026 \textbf{Subject:} Andrei Okounkov (b. July 26, 1969, Moscow, Russia) \textbf{Position:} Professor of Mathematics, Columbia University (since 2010) \textbf{Doctoral Advisor:} Alexander Kirillov (Moscow State University) \textbf{Fields Medal Citation:} \emph{``For his contributions to bridging probability, representation theory and algebraic geometry.''}

\begin{center}\rule{0.5\linewidth}{0.5pt}\end{center}

\begin{prereqs} Representation theory (symmetric groups, infinite-dimensional algebras, Young diagrams), probability theory (random partitions, Plancherel measure), algebraic geometry (Gromov--Witten invariants, moduli spaces). \end{prereqs}

\begin{objectives} Explain the Plancherel measure and its limit shape; describe the Okounkov--Vershik approach to the symmetric group; prove the MNOP conjecture relating Gromov--Witten and Donaldson--Thomas invariants; understand Okounkov bodies and the Nekrasov--Okounkov formula. \end{objectives}

\section{Problem History}

Okounkov's work \cite{okounkov2003} spans three major areas that had previously been considered separate: the asymptotic representation theory of the symmetric groups, enumerative algebraic geometry, and integrable probability. In each area, he uncovered deep connections between discrete combinatorial structures and continuous geometric objects.

The asymptotic representation theory of the symmetric group $S_n$ was initiated by Vershik and Kerov in the 1970s, who discovered that random partitions under the Plancherel measure exhibit a deterministic limit shape as $n \to \infty$. The Plancherel measure assigns probability $(\dim \lambda)^2/n!$ to a partition $\lambda \vdash n$. The limit shape, independently found by Okounkov, Stanley, and Kerov--Vershik, coincides with the semicircle distribution from random matrix theory, but the detailed fluctuations---including the Tracy--Widom distribution for the longest row---remained to be established.

In enumerative geometry, the question of counting rational curves on algebraic varieties, made rigorous through Gromov--Witten theory in the 1990s, produced invariants defined via virtual fundamental classes on moduli spaces of stable maps. Independently, Donaldson--Thomas theory counted ideal sheaves of curves on threefolds. The \emph{MNOP conjecture} (2003) proposed a surprising duality between these two theories, relating curve-counting to sheaf-counting through a simple change of variables.

The Virasoro algebra, the central extension of the Lie algebra of polynomial vector fields on the circle, plays a key role in both conformal field theory and the theory of symmetric functions. The Nekrasov partition function in supersymmetric gauge theory can be expressed as a tau function of the Virasoro constraints, and Okounkov's work revealed that these constraints arise from the geometry of moduli spaces of sheaves on algebraic surfaces.

The \emph{universal enveloping algebra} $U(\mathfrak{g})$ of a Lie algebra $\mathfrak{g}$ plays a central role in Okounkov's work on geometric representation theory. The Yangian $Y(\mathfrak{g})$ and the quantum loop algebra $U_q(L\mathfrak{g})$ are two important deformations of $U(\mathfrak{g})$ that arise in integrable systems and knot theory. Okounkov's geometric construction of R-matrices via stable envelopes provides a unified construction of these algebraic structures from the geometry of moduli spaces of sheaves on surfaces.

The \emph{stable envelope} of a fixed point $F$ in a torus-equivariant K-theory module is a map $\operatorname{Stab}_F$ from the K-theory of the ambient space to the K-theory of the fixed locus. The composition of two stable envelopes factors through the fixed loci, and the resulting composition rules give the R-matrix of the quantum group. This construction, developed by Okounkov and Maulik, is remarkable because it produces purely algebraic objects (R-matrices, quantum group representations) from purely geometric data (moduli spaces, fixed points, intersection theory).

\section{How Was It Solved}

\textbf{Plancherel measure and random partitions.} Okounkov established explicit asymptotics for the correlation functions of the Plancherel measure in terms of the discrete sine kernel. He proved that the distribution of the first row of a random partition follows the Tracy--Widom distribution from random matrix theory, and he identified precise correspondences between fluctuations of random partitions and eigenvalues of random matrices from the Gaussian Unitary Ensemble. The Tracy--Widom distribution, $F_2(s) = \exp(-\int_s^\infty (x-s)q(x)\,dx)$ where $q$ satisfies the Painlev\u00e9 II equation, arises as the limiting distribution of the largest eigenvalue in the GUE; Okounkov showed it governs the longest row of a random partition under the Plancherel measure. The key technical tool is the \emph{transition measure} on the Young diagram, which encodes the branching probabilities from $S_n$ to $S_{n-1}$ and satisfies a discrete integrable system.

The \emph{Okounkov--Vershik approach.} In a 1996 landmark paper with Vershik, Okounkov provided a new construction of the irreducible representations of $S_n$ using the Gelfand--Tsetlin basis and the chain of subgroups $S_1 \subset S_2 \subset \cdots \subset S_n$. He proved that the set of all central measures on the Young graph is in bijection with the Thoma simplex, classifying all possible characters of the infinite symmetric group $S_\infty$. This became the standard modern framework for the subject. The Thoma simplex is parametrized by sequences $(\alpha_1, \alpha_2, \ldots)$ and $(\beta_1, \beta_2, \ldots)$ with $\sum \alpha_i + \sum \beta_i \le 1$, and it classifies all extremal characters of $S_\infty$.

The \emph{Gelfand--Tsetlin basis} for an irreducible representation $V^\lambda$ of $S_n$ is indexed by Young tableaux of shape $\lambda$. The basis vectors $v_T$ are labeled by standard Young tableaux $T$ of shape $\lambda$, and the branching rule $\operatorname{Res}^{S_n}_{S_{n-1}} V^\lambda = \bigoplus_{\mu} V^\mu$ (where the sum is over partitions $\mu$ obtained by removing one box from $\lambda$) determines the action of the Jucys--Murphy elements $J_k = \sum_{j < k} (j,k)$ on the basis vectors. The eigenvalue of $J_k$ on $v_T$ is the content $c(T(k))$ of the box containing $k$ in $T$. This elegant formula, which Okounkov and Vershik derived using the Gelfand--Tsetlin basis, is the foundation for the asymptotic analysis of representations of $S_n$.

The \emph{Young graph} $Y$ is the graded graph whose vertices at level $n$ are partitions $\lambda \vdash n$, with edges given by the inclusion relation $\mu \subset \lambda$ when $\lambda$ is obtained from $\mu$ by adding one box. Central measures on $Y$ are probability measures $\mu$ such that $\mu(\lambda) \dim \lambda / n!$ is a class function on $S_n$ for each $n$. The classification of central measures on $Y$ as convex combinations of extremal measures (which correspond to points of the Thoma simplex) is the \emph{Thoma theorem}, and Okounkov's contribution was to provide a new proof using the Gelfand--Tsetlin basis and to extend the result to asymptotic characters.

\textbf{The MNOP conjecture.} In 2003, Maulik, Nekrasov, Okounkov, and Pandharipande conjectured that for a nonsingular projective threefold $X$, the generating functions of Gromov--Witten invariants $Z_{GW}(X; q)$ and Donaldson--Thomas invariants $Z_{DT}(X; q)$ satisfy:
\[
Z_{GW}(X; q) = Z_{DT}(X; -q)^{\chi(X)}
\]
where $\chi(X)$ is the topological Euler characteristic. Okounkov and Pandharipande proved this for local curves, toric threefolds, and Calabi--Yau threefolds using virtual localization on moduli spaces of stable pairs, combined with sophisticated combinatorial analysis of fixed point contributions. The proof uses equivariant localization with respect to the natural torus action, reducing the virtual integrals to combinatorial sums over plane partitions (3D Young diagrams) that index fixed points.

\textbf{The Nekrasov--Okounkov formula.} With Nekrasov, Okounkov proved that the instanton partition function of $N=2$ supersymmetric Yang--Mills theory, in the limit $\epsilon_1, \epsilon_2 \to 0$, reproduces the Seiberg--Witten prepotential. The proof uses equivariant localization to reduce integrals to sums over pairs of partitions, and the Cauchy identity to simplify the expression to the Seiberg--Witten form. The Nekrasov--Okounkov formula expresses the instanton partition function as:
\[
Z_{\text{inst}} = \sum_{\vec{Y}} \prod_{(i,j) \in Y} \frac{1}{a(Y_{i,j}) \cdot (a(Y_{i,j}) + \epsilon_1 + \epsilon_2)},
\]
where the sum is over tuples of partitions $\vec{Y} = (Y_1, \ldots, Y_N)$ and $a(i,j)$ denotes the arm and leg lengths of the box $(i,j)$.

\textbf{Okounkov bodies.} In 2003, Okounkov introduced convex bodies $\Delta(L) \subset \mathbb{R}^n$ associated to line bundles $L$ on projective varieties, whose Euclidean volume equals the volume of $L$. These \emph{Okounkov bodies} provide a convex-geometric approach to volumes, positivity, and asymptotic invariants, with applications to Fujita's approximation theorem and Zariski decompositions. For a line bundle $L$ on a surface $X$, the Okounkov body $\Delta(L) \subset \mathbb{R}^2$ is constructed by choosing a valuation on the function field and collecting all valuation vectors of sections of $L$. The key property is that the lattice points in $\Delta(L)$ correspond to sections of $L$, and the volume of $\Delta(L)$ equals the self-intersection number $L^2/2$. This construction generalizes the classical Newton polytope of a polynomial and provides a unified framework for studying linear series on varieties of arbitrary dimension.

The construction of Okounkov bodies proceeds as follows. Fix a flag of subvarieties $X = X_d \supset X_{d-1} \supset \cdots \supset X_0 = \{pt\}$ in $X$. For each section $s \in H^0(X, L)$, define the valuation vector $v(s) = (v_1(s), \ldots, v_d(s)) \in \mathbb{R}^d$, where $v_i(s)$ is the order of vanishing of $s$ along $X_i$ (relative to $X_{i-1}$). The Okounkov body is then the convex hull of all valuation vectors:
\[
\Delta(L) = \overline{\bigcup_{s \in H^0(X,L) \setminus \{0\}} \frac{v(s)}{m} + [0,1)^d},
\]
where the closure is taken in $\mathbb{R}^d$. The key theorem is that the Euclidean volume of $\Delta(L)$ equals the volume of $L$ normalized by $d!$, and that $\Delta(L)$ is independent of the choice of flag for nef line bundles.

\textbf{Gromov--Witten invariants and stable maps.} Okounkov's work on Gromov--Witten theory concerns the enumeration of maps from smooth curves to a target variety $X$. The moduli space $\overline{\mathcal{M}}_{g,n}(X, \beta)$ of stable maps of genus $g$ with $n$ marked points representing class $\beta \in H_2(X, \mathbb{Z})$ carries a virtual fundamental class of the expected dimension, and the Gromov--Witten invariants are the intersection numbers on this virtual class. Okounkov's contribution was to connect these invariants to the representation theory of symmetric groups and to the combinatorics of partitions, revealing unexpected links between enumerative geometry and algebraic combinatorics.

\section{Impact of This Discovery}

The MNOP conjecture and its proof transformed enumerative geometry by establishing Donaldson--Thomas theory as a fundamental tool and revealing a deep unity between curve-counting and sheaf-counting. The techniques---virtual localization, the Fock space approach, and the Virasoro constraints---have become standard in the field.

The Okounkov--Vershik approach to the symmetric group provided the definitive modern framework, clarifying the role of the Young--Jucys--Murphy elements and opening the door to asymptotic theory. The geometric construction of R-matrices using stable envelopes, developed in later work with Maulik, established a bridge between enumerative geometry and the representation theory of quantum groups. Okounkov's work on random partitions showed that the Plancherel measure satisfies a determinantal structure connecting it to random matrix theory, and his construction of Newton--Okounkov bodies provided a convex-geometric approach to volumes of line bundles on algebraic varieties.

The \emph{Virasoro constraints} are an infinite system of differential equations satisfied by the partition function of a conformal field theory. The Nekrasov partition function in supersymmetric gauge theory satisfies Virasoro constraints, and Okounkov's work showed that these constraints arise from the geometry of the moduli space of sheaves on algebraic surfaces. The Virasoro algebra is the central extension of the Lie algebra of polynomial vector fields on the circle, and its representation theory is fundamental to two-dimensional conformal field theory. Okounkov's work connects the Virasoro constraints to the Donaldson--Thomas invariants via the MNOP correspondence, and the Virasoro constraints provide a powerful tool for computing the partition function and its derivatives.

The \emph{Fock space representation} of the Heisenberg algebra is the algebraic structure underlying the computation of Gromov--Witten invariants via localization. The Fock space is the symmetric algebra $\operatorname{Sym}(H_2(X, \mathbb{Z}))$, and the Heisenberg generators correspond to the insertion operators in the Gromov--Witten theory. Okounkov and Pandharipande showed that the partition function of the Gromov--Witten theory can be expressed as a vacuum expectation value in the Fock space, and the Virasoro constraints are realized as commutation relations in the Fock space representation. This algebraic structure provides a powerful computational tool for enumerative geometry.

Okounkov's work on random partitions and the Plancherel measure connected probability theory with algebraic combinatorics, making the Tracy--Widom distribution, the limit shape, and the frozen boundary of random tilings central objects in integrable probability. The Nekrasov--Okounkov formula provided rigorous mathematical foundations for instanton counting in supersymmetric gauge theory, with applications to the AGT correspondence and topological string theory.

The \emph{Plancherel measure} on partitions of $n$ is the probability measure that assigns to each partition $\lambda \vdash n$ the weight $(\dim \lambda)^2 / n!$, where $\dim \lambda$ is the number of standard Young tableaux of shape $\lambda$. The Plancherel measure arises naturally in representation theory as the distribution of the irreducible representation of $S_n$ that appears with highest multiplicity when a random tensor product of representations is decomposed. Okounkov showed that the Plancherel measure satisfies a determinantal structure: the correlation functions of the Plancherel measure are determinants of the discrete sine kernel $\sin(\pi(x-y))/(\pi(x-y))$, which is the same kernel that appears in random matrix theory. This determinantal structure is the source of the integrability of the Plancherel measure and connects the asymptotic representation theory of symmetric groups to the theory of random matrices.

The \emph{discrete sine kernel} $K(x,y) = \sin(\pi(x-y))/(\pi(x-y))$ is the correlation kernel for the eigenvalues of large random matrices from the Gaussian Unitary Ensemble (GUE). Okounkov showed that the same kernel appears in the correlation functions of the Plancherel measure: the probability that two boxes $(i_1, j_1)$ and $(i_2, j_2)$ in a random Young diagram are both occupied is given by a determinant of the discrete sine kernel evaluated at the contents of the boxes. This connection between the Plancherel measure and the GUE is the basis for the Tracy--Widom distribution for the longest row of a random partition.

The \emph{limit shape} of random tilings---the deterministic shape that a random tiling approaches as the domain grows---is described by an algebraic curve in the plane. The frozen boundary, where the tiling transitions from a rigid (frozen) region to a liquid region, is a real algebraic curve whose equation encodes the statistical mechanics of the tiling. Kenyon, Okounkov, and Sheffield showed that the frozen boundary of the dimer model on the hexagonal lattice is described by a certain elliptic curve, and that the fluctuations near the frozen boundary are governed by the Tracy--Widom distribution from random matrix theory. The frozen boundary has a beautiful geometric description: it is the set of points where the height function of the tiling is not differentiable, and it separates the ``arctic circle'' (the frozen region) from the ``liquid region'' (where the tiling fluctuates randomly).

The \emph{Schur process} is a probability measure on sequences of partitions that encodes the correlation functions of the Plancherel measure and other random partition ensembles. Okounkov showed that the Schur process satisfies a determinantal structure, meaning that the correlation functions are determinants of a certain kernel matrix. This determinantal structure is the source of the integrability of the Plancherel measure, and it connects the asymptotic representation theory of symmetric groups to the theory of orthogonal polynomials and random matrices. The \emph{AGT correspondence} (Alba--Gaiotto--Tachikawa, 2010) relates the Nekrasov partition function of $N=2$ supersymmetric Yang--Mills theory on $\mathbb{C}^2$ to the conformal blocks of a two-dimensional conformal field theory. This correspondence, which was first observed by physicists and later proved in special cases, connects the four-dimensional gauge theory to a two-dimensional CFT via the geometry of class $\mathcal{S}$ theories obtained by compactifying six-dimensional $\mathcal{N}=(2,0)$ theory on a Riemann surface.

The \emph{topological string amplitude} on a Calabi--Yau threefold $X$ is a generating function of Gromov--Witten invariants (or equivalently, of Donaldson--Thomas invariants via the MNOP correspondence). Okounkov's work on the MNOP conjecture provided the mathematical framework for understanding the topological string amplitude, and the relationship between topological strings and supersymmetric gauge theory (via geometric engineering) connects his work to string theory and quantum field theory.

The Okounkov body construction has become a fundamental tool in birational geometry, used to prove the Khovanskii--Teissier conjecture on concavity of intersection numbers and to study mirror symmetry and toric degenerations. The Newton--Okounkov body generalization of classical Newton polytopes has found applications throughout algebraic geometry.

\section{Modern Views and Improvements}

Okounkov continues to develop the theory of stable envelopes and their applications to geometric representation theory. The elliptic stable envelopes, a generalization to the elliptic setting, have applications to elliptic quantum groups and 3D mirror symmetry. The geometric R-matrix construction using stable envelopes produces R-matrices for Yangians, quantum loop algebras, and elliptic quantum groups of $\mathfrak{gl}_n$.

The frozen boundary of random tilings, studied with Kenyon and Sheffield, connects algebraic geometry, complex analysis, and probability theory. The limit shape of random tilings and the description of the frozen boundary by algebraic curves are now central objects in the study of random interfaces.

Okounkov bodies have become a fundamental tool in birational geometry, used to prove the Khovanskii--Teissier conjecture on concavity of intersection numbers and to study mirror symmetry and toric degenerations. The Newton--Okounkov body generalization of classical Newton polytopes has found applications throughout algebraic geometry. Recent work by K\u00f6nig, L\"ubke, and others extends the theory to non-abelian valuations and higher-dimensional varieties.

The \emph{Khovanskii--Teissier conjecture} states that for a nef line bundle $L$ on a projective variety $X$ of dimension $n$, the function $k \mapsto (L^k)^{1/k}$ is log-concave. Okounkov bodies provide a proof of this conjecture: the volume of $\Delta(kL) = k^n \Delta(L)$ is $k^n \operatorname{Vol}(\Delta(L))$, and the log-concavity follows from the Brunn--Minkowski inequality for convex bodies. The key insight is that the Okounkov body $\Delta(L)$ captures the essential geometric information about $L$, and the volume of $\Delta(L)$ equals the self-intersection number $\deg(L)/n!$. This approach, which Okounkov developed in his 2003 paper, provides a unified framework for studying linear series on varieties of arbitrary dimension.

The \emph{Newton--Okounkov body} generalization of classical Newton polytopes has applications throughout algebraic geometry. For a polynomial $f$ in $n$ variables, the Newton polytope is the convex hull of the exponent vectors of the monomials appearing in $f$. The Okounkov body of the associated line bundle on $\mathbb{P}^n$ is a translated copy of the Newton polytope, and the lattice points in the Okounkov body correspond to the monomials appearing in sections of the line bundle. This connection between convex geometry and algebraic geometry has applications to mirror symmetry, toric degenerations, and the study of K\u00e4hler--Einstein metrics.

The \emph{AGT correspondence} (Alba--Gaiotto--Tachikawa, 2010) relates the Nekrasov partition function of $N=2$ supersymmetric Yang--Mills theory on $\mathbb{C}^2$ to the conformal blocks of a two-dimensional conformal field theory. This correspondence, which was first observed by physicists and later proved in special cases, connects the four-dimensional gauge theory to a two-dimensional CFT via the geometry of class $\mathcal{S}$ theories obtained by compactifying six-dimensional $\mathcal{N}=(2,0)$ theory on a Riemann surface. Okounkov's work on the Nekrasov--Okounkov formula provides the mathematical framework for the AGT correspondence, and the connection between instanton counting and conformal blocks is an active area of research.

Recent progress by Okounkov and collaborators on \emph{3D mirror symmetry} and \emph{cohomological Hall algebras} continues to bridge representation theory, algebraic geometry, and mathematical physics. The cohomological Hall algebra of a Calabi--Yau 3-category provides a categorical framework for the Donaldson--Thomas invariants and the BPS state counting.

The \emph{BPS state counting} problem, which arises in string theory, asks for the number of BPS states with given charge in a compactified string theory. Okounkov's work on Donaldson--Thomas invariants provides a mathematical framework for this counting, and the MNOP conjecture relates it to the Gromov--Witten invariants that arise from topological string theory. The partition function of BPS states is expected to transform as a modular form under the duality group, and Okounkov's work provides evidence for this modularity in specific examples.

The \emph{plane partition} $\pi$ is an array of non-negative integers $\pi_{i,j}$ that is non-increasing in both rows and columns. The generating function of plane partitions, proved by MacMahon, is $\sum_\pi q^{|\pi|} = \prod_{n=1}^\infty (1-q^n)^{-n}$, and this product formula is the starting point for the equivariant localization computation in the MNOP conjecture. The fixed points of the torus action on the moduli space of ideal sheaves are indexed by plane partitions, and the weights of the fixed points are expressed in terms of the arm and leg lengths of the boxes in the 3D Young diagram. The key identity in Okounkov and Pandharipande's proof is the evaluation of the weighted sum over plane partitions as a product formula, which they establish using a residue calculus argument and the Cauchy determinant identity.

The \emph{quantum cohomology} of a smooth variety $X$ is a deformation of the classical cohomology ring using Gromov--Witten invariants as structure constants: $\alpha * \beta = \sum_\beta \langle \alpha, \beta, \gamma \rangle_\beta \gamma^*$, where the sum is over curve classes and $\gamma^*$ is the Poincar\'{e} dual basis. Okounkov's work connects quantum cohomology to the representation theory of quantum groups via the quantum Riemann--Roch theorem and the theory of stable envelopes. The quantum cohomology of a flag variety $G/B$ is isomorphic to the center of the appropriate quantum group, and this connection provides a representation-theoretic interpretation of the enumerative invariants.

The \emph{random surface} is a random planar map (a graph embedded in the sphere) considered as a metric space using the graph distance. The Brownian map, introduced by Le Gall and Miermont, is the scaling limit of uniformly random planar maps. Okounkov's work on the Nekrasov partition function provides a connection between random surfaces and representation theory: the partition function of the random surface can be expressed as a sum over Young diagrams, and the fluctuations of the surface are governed by the Plancherel measure and the Tracy--Widom distribution. This connection, developed in collaboration with Borodin and Corwin, has applications to integrable probability and to the study of random matrices.

The \emph{integrable hierarchy} associated to the KP (Kadomtsev--Petviashvili) equation is a sequence of commuting differential equations whose solutions are $\tau$-functions. The Schur process, introduced by Okounkov and Reshetikhin, is a probability measure on sequences of Young diagrams that satisfies the KP hierarchy. The key property is that the correlation functions of the Schur process are determinants of a certain kernel, and this determinantal structure is the source of the integrability. The Schur process provides a unified framework for studying random partitions, random matrices, and the Plancherel measure, and it connects Okounkov's work on random partitions to the theory of integrable systems.

\begin{center}\rule{0.5\linewidth}{0.5pt}\end{center}

\section*{Summary}
\addcontentsline{toc}{section}{Summary}

Okounkov established deep connections between representation theory, probability, and algebraic geometry through his work on random partitions, the MNOP conjecture relating Gromov--Witten and Donaldson--Thomas invariants, and the Nekrasov--Okounkov formula for instanton counting. His contributions continue to shape enumerative geometry, integrable probability, and geometric representation theory.
